\documentclass[conference]{IEEEtran}
\usepackage{cite}
\usepackage{amsmath,amssymb,amsfonts}
\usepackage{algorithmic}
\usepackage{graphicx}
\usepackage{textcomp}
\usepackage{xcolor}
\usepackage{booktabs} % For nicer tables
\usepackage{float}

\begin{document}

% Title remains the same
\title{DynaLITE: A miniature DynamoDB with replication}

\author{
    % All names in a single line separated by commas
    \IEEEauthorblockN{Kushagra Trivedi, Ronit Jalihal, and Prasoon Dev}
    % All affiliations/student IDs in the next block, separated by spaces or line breaks
    \IEEEauthorblockA{\textit{2023101048 \qquad 2023101028 \qquad 2023111014}}
    % Separate line break for the supervisor/project info
    \\[0.5ex]
    \parbox{\linewidth}{\centering
        \vspace{0.2cm}
        \footnotesize
        Distributed System project under Dr. Kishore Kothapalli
    }
}

\maketitle

\begin{abstract}
DynamoMINI is an educational implementation of Amazon's Dynamo, a highly available key-value storage system designed to provide scalability, fault tolerance, and low latency. The system demonstrates core distributed system principles including consistent hashing for data partitioning, vector clocks for version tracking, quorum-based replication, gossip protocols for membership management, hinted handoff for handling temporary failures, and Merkle trees for anti-entropy. This comprehensive report presents the architectural design, implementation details, and operational characteristics of DynamoMINI. We explore the fundamental trade-offs between consistency, availability, and partition tolerance (CAP theorem), demonstrating how DynamoDB prioritizes availability and partition tolerance while achieving eventual consistency. The report includes detailed system design diagrams illustrating data flow, replication strategies, failure handling mechanisms, and background synchronization processes. Through this implementation, we validate the feasibility of building a production-grade distributed storage system that can handle node failures, network partitions, and dynamic scaling while maintaining high availability for both read and write operations.
\end{abstract}

\section{Introduction}

\subsection{Background and Motivation}
Amazon's Dynamo represents a paradigm shift in distributed database design, prioritizing availability and partition tolerance over strict consistency. Introduced to power Amazon's critical services including the shopping cart, session management, and product catalog, Dynamo addresses the unique requirements of large-scale e-commerce platforms where service availability directly impacts revenue. Traditional RDBMS solutions, while providing ACID guarantees, often sacrifice availability during network partitions or node failures—an unacceptable trade-off for services requiring ``always-on'' operation.

DynamoDB's design philosophy stems from the realization that many applications can tolerate eventual consistency in exchange for higher availability and lower latency. For instance, in a shopping cart scenario, it is preferable to accept a slightly stale version of the cart than to reject a user's request to add items during a network partition. This approach embraces the CAP theorem, explicitly choosing availability and partition tolerance while accepting eventual consistency.

The system employs sophisticated techniques to manage the complexity introduced by this design choice: consistent hashing for incremental scalability, vector clocks for tracking causality and detecting conflicts, quorum-based replication for tunable consistency, gossip protocols for decentralized coordination, hinted handoff for handling temporary failures, and Merkle trees for efficient anti-entropy repair. These mechanisms work in concert to provide a robust, self-healing distributed storage system.

\subsection{DynamoMINI: An Educational Implementation}
DynamoMINI is a fully functional implementation of Amazon's Dynamo, designed to demonstrate and validate the core principles of highly available distributed storage systems. Built using Python, RPyC for remote procedure calls, and Redis as the underlying storage engine, DynamoMINI implements the complete architecture including:

\begin{itemize}
    \item \textbf{Consistent Hashing Ring:} Complete implementation with virtual nodes for load balancing
    \item \textbf{Quorum-based Replication:} Configurable N, R, W parameters for tunable consistency
    \item \textbf{Vector Clock Versioning:} Full causality tracking and conflict detection
    \item \textbf{Gossip Protocol:} Decentralized membership and failure detection
    \item \textbf{Hinted Handoff:} Temporary failure handling mechanism
    \item \textbf{Merkle Tree Anti-Entropy:} Background synchronization for replica consistency
    \item \textbf{Background Threads:} Asynchronous processes for system maintenance
\end{itemize}

This implementation serves both as an educational tool for understanding distributed systems principles and as a validation of the feasibility of building production-grade distributed storage systems. Through DynamoMINI, we explore the practical challenges of implementing theoretical concepts, the trade-offs inherent in distributed system design, and the operational characteristics of eventually consistent storage systems.

% Figure: Redundancy Diagram
\begin{figure}[htbp]
    \centering
    \includegraphics[width=\columnwidth]{redundant_diagram.png}
    \caption{DynamoMINI redundancy and replication architecture showing how data is replicated across multiple nodes for fault tolerance and high availability.}
    \label{fig:redundancy}
\end{figure}

\section{Theoretical Foundations}

\subsection{The CAP Theorem and Design Trade-offs}
The CAP theorem, formulated by Eric Brewer, states that a distributed system can provide at most two of three guarantees: Consistency, Availability, and Partition tolerance. DynamoDB explicitly chooses availability and partition tolerance, accepting eventual consistency as the trade-off. This decision is driven by the operational requirements of large-scale web services where:

\begin{itemize}
    \item Service availability directly impacts user experience and revenue
    \item Network partitions are inevitable in geographically distributed systems
    \item Many applications can tolerate temporary inconsistencies
    \item Conflict resolution can be pushed to read operations or application logic
\end{itemize}

\subsection{Consistent Hashing}
Consistent hashing solves the fundamental problem of data distribution in dynamic distributed systems. Traditional modulo-based hashing requires complete data redistribution when the number of nodes changes. Consistent hashing minimizes this overhead by mapping both data keys and nodes onto a circular hash space.

\textbf{Key Properties:}
\begin{itemize}
    \item \textbf{Load Distribution:} Keys are uniformly distributed across nodes
    \item \textbf{Minimal Disruption:} Adding/removing a node affects only immediate neighbors
    \item \textbf{Virtual Nodes:} Each physical node is assigned multiple positions on the ring to improve load balancing
    \item \textbf{Incremental Scalability:} The system scales linearly with the number of nodes
\end{itemize}

In DynamoMINI, MD5 hashing maps keys to a 128-bit hash space. Each node claims multiple virtual positions on the ring, ensuring that even with a small number of physical nodes, the load distribution remains balanced. When a key is accessed, its hash determines the coordinator node—the first node encountered by moving clockwise from the key's position.

% Figure: General Hashing Sequence
\begin{figure}[htbp]
    \centering
    \includegraphics[width=0.85\columnwidth]{hashing.png}
    \caption{General hashing sequence showing how keys are mapped to hash values and subsequently to nodes in the system.}
    \label{fig:hashing_sequence}
\end{figure}

% Figure 1: Consistent Hashing Ring
\begin{figure}[htbp]
    \centering
    \includegraphics[width=0.9\columnwidth]{fig1.png}
    \caption{Consistent hashing ring showing data partitioning and replication. Nodes are positioned via MD5 hash, keys are assigned to the first N nodes clockwise from the key's hash position.}
    \label{fig:consistent_hashing}
\end{figure}

\subsection{Vector Clocks for Causality Tracking}
Vector clocks provide a mechanism for tracking causality and detecting conflicts in distributed systems. Unlike physical timestamps, vector clocks capture the happens-before relationship between events across different nodes.

\textbf{Vector Clock Structure:} Each data version is associated with a vector clock: $VC = [N_1: v_1, N_2: v_2, ..., N_k: v_k]$ where $N_i$ represents a node and $v_i$ is the version counter for that node.

\textbf{Operations:}
\begin{itemize}
    \item \textbf{Update:} When node $N_i$ modifies data, it increments $v_i$
    \item \textbf{Merge:} When receiving data from node $N_j$, node $N_i$ takes the maximum of each counter
    \item \textbf{Comparison:} $VC_1 < VC_2$ if all counters in $VC_1$ are $\leq$ corresponding counters in $VC_2$, with at least one strictly less
    \item \textbf{Conflict Detection:} If neither $VC_1 < VC_2$ nor $VC_2 < VC_1$, versions are concurrent (conflicting)
\end{itemize}

% Figure 2: Vector Clock Evolution
\begin{figure}[htbp]
    \centering
    \includegraphics[width=0.85\columnwidth]{fig2.png}
    \caption{Version evolution of an object showing vector clock progression. Concurrent writes create version branches that must be reconciled during reads.}
    \label{fig:vector_clocks}
\end{figure}

\subsection{Quorum-Based Replication}
Quorum protocols provide tunable consistency guarantees by requiring a minimum number of nodes to participate in read and write operations. The system uses three parameters:

\begin{itemize}
    \item $N$: Total number of replicas (replication factor)
    \item $W$: Write quorum—minimum nodes that must acknowledge a write
    \item $R$: Read quorum—minimum nodes that must respond to a read
\end{itemize}

\textbf{Consistency Guarantee:} By ensuring $R + W > N$, the system guarantees that any read quorum overlaps with any write quorum, ensuring that reads see the most recent writes.

\textbf{Configuration Examples:}
\begin{itemize}
    \item $R=1, W=N$: Fast reads, slow writes, optimized for read-heavy workloads
    \item $R=N, W=1$: Fast writes, slow reads, optimized for write-heavy workloads
    \item $R=W=\lceil(N+1)/2\rceil$: Balanced approach with equal latency for reads/writes
\end{itemize}

% Figure 3: Quorum Technique
\begin{figure}[htbp]
    \centering
    \includegraphics[width=0.7\columnwidth]{quorum.png}
    \caption{Quorum technique with N=3, R=2, W=2. The overlap between read and write quorums guarantees consistency.}
    \label{fig:quorum}
\end{figure}

\section{Basic Assumptions}
DynamoMINI is designed as a distributed key-value store, where data is partitioned and replicated across multiple nodes to ensure fault tolerance and high availability. It operates under the following assumptions: % [cite: 63, 64]
\begin{itemize}
    \item \textbf{Unique Identification:} Each data item is uniquely identified by a key. % [cite: 65]
    \item \textbf{No Cross-Item Operations:} Operations are confined to single data items. % [cite: 66]
    \item \textbf{Trusted Environment:} Security concerns are out of scope; assume a trusted environment. % [cite: 67]
\end{itemize}

\section{System Architecture}

\subsection{Overall System Design}
DynamoMINI follows a fully decentralized, peer-to-peer architecture where each node has identical responsibilities and capabilities. This symmetric design eliminates single points of failure and simplifies operational management. The architecture comprises several interconnected subsystems working in concert to provide high availability and eventual consistency.

% Figure 4: Overall Architecture
\begin{figure}[htbp]
    \centering
    \includegraphics[width=\columnwidth]{architecture.png}
    \caption{DynamoMINI system architecture showing decentralized node structure with internal components and inter-node communication patterns.}
    \label{fig:architecture}
\end{figure}

\subsection{Data Partitioning and Replication}
Data is partitioned across nodes using consistent hashing, with each key assigned to a \textit{preference list}—the N successive nodes encountered clockwise from the key's hash position. The first node in the preference list acts as the coordinator for that key, handling client requests and orchestrating replication.

\textbf{Preference List Construction:}
\begin{enumerate}
    \item Hash the key using MD5 to determine position on ring
    \item Starting from that position, traverse clockwise collecting N distinct physical nodes
    \item Skip virtual nodes belonging to already-selected physical nodes
    \item The resulting list becomes the preference list for that key range
\end{enumerate}

% Figure 5: Data Partitioning
\begin{figure}[htbp]
    \centering
    \includegraphics[width=0.9\columnwidth]{fig3.png}
    \caption{Data partitioning and replication on the hash ring. Key K's preference list consists of nodes A, B, C (N=3), which store replicas of all keys in the shaded range.}
    \label{fig:partitioning}
\end{figure}

\subsection{Request Routing}
Clients can contact any node in the system. If the contacted node is not in the key's preference list, it forwards the request to the appropriate coordinator. This approach provides load balancing while maintaining efficient request routing.

\textbf{Routing Decision Process:}
\begin{enumerate}
    \item Receive client request for key K
    \item Compute hash(K) and identify preference list from local routing table
    \item If self is coordinator (first in preference list), handle request locally
    \item Otherwise, forward request to coordinator node
    \item Coordinator executes request and returns response to client
\end{enumerate}

\subsection{Replication Strategy}
DynamoMINI implements asynchronous replication with configurable quorum parameters. Write operations proceed in two phases:

\textbf{Synchronous Phase:} The coordinator writes locally and sends replication requests to N-1 replicas, waiting for W-1 acknowledgments before responding to the client.

\textbf{Asynchronous Phase:} Remaining replicas are updated in the background. If nodes are temporarily unavailable, hinted handoff queues the writes for later delivery.

% Figure 6: Master System Flow
\begin{figure}[htbp]
    \centering
    \includegraphics[width=\columnwidth]{system_diagram_events.png}
    \caption{Complete system design integrating all components: request routing, consistent hashing, quorum logic, and background protocols.}
    \label{fig:system_flow}
\end{figure}

\subsection{Gossip Protocol for Membership Management}
The gossip protocol maintains a consistent view of cluster membership without centralized coordination. Each node periodically:

\begin{enumerate}
    \item Selects a random peer from its routing table
    \item Exchanges routing table information including node list and version vectors
    \item Compares received information with local state
    \item Updates local routing table if peer has newer information
    \item Merges down node lists to track potentially failed nodes
\end{enumerate}

\textbf{Convergence Properties:} The gossip protocol exhibits logarithmic convergence—information propagates to all nodes in $O(\log N)$ gossip rounds, where N is the number of nodes. This rapid dissemination ensures that topology changes are quickly reflected across the cluster.

% Figure 7: Gossip Protocol
\begin{figure}[htbp]
    \centering
    \includegraphics[width=0.85\columnwidth]{vector_clocks.png}
    \caption{Gossip protocol execution flow showing periodic information exchange and routing table reconciliation.}
    \label{fig:gossip}
\end{figure}

\subsection{Failure Detection and Recovery}
DynamoMINI employs a multi-layered approach to failure detection:

\textbf{Passive Detection:} Failed RPC calls during normal operations mark nodes as potentially down.

\textbf{Active Probing:} A dedicated background thread periodically pings nodes in the ``down'' list, attempting to detect recovery.

\textbf{Gossip-Based Dissemination:} Failure information propagates via gossip, ensuring all nodes maintain consistent failure views.

\textbf{Hinted Handoff:} During temporary failures, writes intended for unavailable nodes are stored as ``hints'' on alternative nodes and replayed upon recovery.

\subsection{Background Synchronization}
Three background threads maintain system health:

\textbf{Gossip Thread:} Runs every few seconds, exchanges membership information with random peers.

\textbf{Anti-Entropy Thread:} Periodically builds Merkle trees over key ranges and exchanges tree summaries with neighbors to identify and repair inconsistencies.

\textbf{Failure Detection Thread:} Continuously monitors the down-node list, attempting reconnection and triggering hinted handoff replay upon success.

% Figure 8: Background Threads
\begin{figure}[htbp]
    \centering
    \includegraphics[width=\columnwidth]{async_threads.png}
    \caption{State machines for three background threads ensuring system maintenance and consistency.}
    \label{fig:background_threads}
\end{figure}

\section{Detailed Operation Flows}

\subsection{PUT Operation Execution}
The PUT operation demonstrates how DynamoMINI achieves high write availability while maintaining eventual consistency through quorum-based replication and vector clock versioning.

\textbf{Execution Phases:}

\begin{enumerate}
    \item \textbf{Request Reception:} Client issues PUT(key, value) to any node
    \item \textbf{Coordinator Selection:} Request is forwarded to the first node in the key's preference list if not already there
    \item \textbf{Version Generation:} Coordinator generates new vector clock by:
    \begin{itemize}
        \item Reading existing vector clock(s) for the key
        \item Incrementing its own counter in the vector
        \item Creating new version with updated vector clock
    \end{itemize}
    \item \textbf{Local Write:} Coordinator writes (key, value, vector\_clock, timestamp) to local Redis storage
    \item \textbf{Parallel Replication:} Coordinator sends asynchronous replication requests to N-1 replica nodes
    \item \textbf{Quorum Wait:} Coordinator waits for W-1 acknowledgments (having already written locally, total writes = W)
    \item \textbf{Response:} Upon achieving write quorum, coordinator responds SUCCESS to client
    \item \textbf{Background Updates:} Remaining replicas continue updating asynchronously
    \item \textbf{Failure Handling:} If any replica is unreachable, hint is stored for later replay via hinted handoff
\end{enumerate}

% Figure 9: PUT Sequence Diagram
\begin{figure}[htbp]
    \centering
    \includegraphics[width=0.9\columnwidth]{put_operation_sequence.png}
    \caption{PUT operation sequence showing quorum-based replication with W=2. Coordinator writes locally and waits for one additional replica acknowledgment before responding.}
    \label{fig:put_sequence}
\end{figure}

\subsection{GET Operation Execution}
The GET operation prioritizes returning the most recent version of data while performing background consistency repair (read repair).

\textbf{Execution Phases:}

\begin{enumerate}
    \item \textbf{Request Reception:} Client issues GET(key) to any node
    \item \textbf{Coordinator Selection:} Request forwarded to coordinator (first in preference list)
    \item \textbf{Local Read:} Coordinator reads its local version of the key
    \item \textbf{Parallel Reads:} Coordinator sends asynchronous read requests to N-1 replica nodes
    \item \textbf{Quorum Wait:} Coordinator waits for R responses (including its own local read)
    \item \textbf{Version Reconciliation:} 
    \begin{itemize}
        \item Compare vector clocks from all responding replicas
        \item Identify the most recent version(s) using vector clock comparison
        \item If a single version dominates all others, return it
        \item If multiple concurrent versions exist (conflict), return all conflicting versions for application-level reconciliation
    \end{itemize}
    \item \textbf{Read Repair:} Coordinator identifies replicas with stale versions and asynchronously updates them with the latest version
    \item \textbf{Response:} Return value(s) to client
\end{enumerate}

% Figure 10: GET Sequence Diagram
\begin{figure}[htbp]
    \centering
    \includegraphics[width=0.9\columnwidth]{get_operation_sequence.png}
    \caption{GET operation sequence showing quorum-based read with R=2 and read repair mechanism for consistency maintenance.}
    \label{fig:get_sequence}
\end{figure}

\subsection{Hinted Handoff Mechanism}
Hinted handoff ensures write availability during temporary node failures by temporarily storing writes on alternative nodes and replaying them upon node recovery.

\textbf{Operational Flow:}

\begin{enumerate}
    \item \textbf{Failure Detection:} During PUT replication, coordinator detects that replica node A is unreachable
    \item \textbf{Hint Storage:} Coordinator stores a hint in local Redis: \texttt{hint:nodeA:key → (value, vector\_clock, timestamp)}
    \item \textbf{Hint Logging:} Hint is logged in a persistent queue for the target node
    \item \textbf{Recovery Detection:} Background failure detection thread periodically pings down nodes
    \item \textbf{Node Recovery:} When node A becomes reachable again, coordinator detects recovery
    \item \textbf{Hint Replay:} All hints for node A are retrieved and replayed:
    \begin{itemize}
        \item Send PUT(key, value, vector\_clock) to recovered node
        \item Wait for acknowledgment
        \item Delete hint from local storage upon successful replay
    \end{itemize}
    \item \textbf{Retry Logic:} If replay fails, hint remains queued for subsequent retry
\end{enumerate}

% Figure 11: Hinted Handoff
\begin{figure}[htbp]
    \centering
    \includegraphics[width=0.85\columnwidth]{hinted_handoff.png}
    \caption{Hinted handoff mechanism for handling temporary node failures and ensuring eventual consistency.}
    \label{fig:hinted_handoff}
\end{figure}

\subsection{Merkle Tree Anti-Entropy}
Merkle trees provide an efficient mechanism for detecting and repairing inconsistencies between replicas by comparing cryptographic tree summaries rather than individual keys.

\textbf{Structure and Operation:}

\begin{enumerate}
    \item \textbf{Tree Construction:} Each node periodically builds a Merkle tree over its key range:
    \begin{itemize}
        \item \textbf{Leaf Nodes:} Hash of (value || vector\_clock || timestamp) for each key
        \item \textbf{Internal Nodes:} Hash of concatenated child hashes
        \item \textbf{Root Hash:} Summary of entire key range
    \end{itemize}
    \item \textbf{Tree Exchange:} Nodes exchange Merkle tree roots with their neighbors
    \item \textbf{Difference Detection:} If root hashes differ:
    \begin{itemize}
        \item Recursively compare child node hashes
        \item Identify specific key ranges with inconsistencies
        \item Descend to leaf level to find divergent keys
    \end{itemize}
    \item \textbf{Synchronization:} Exchange and reconcile only the keys with mismatched hashes
    \item \textbf{Efficiency:} For K keys and tree depth D, only $O(\log K)$ hashes need comparison if trees are mostly consistent
\end{enumerate}

% Figure 12: Merkle Tree
\begin{figure}[htbp]
    \centering
    \includegraphics[width=0.9\columnwidth]{merkle_tress.png}
    \caption{Merkle tree-based anti-entropy for efficient replica synchronization. Only divergent branches require data exchange.}
    \label{fig:merkle_tree}
\end{figure}

\subsection{Implementation Technology Stack}

\textbf{Core Components:}
\begin{itemize}
    \item \textbf{Python 3.8+:} Primary implementation language
    \item \textbf{RPyC (Remote Python Call):} RPC framework for inter-node communication
    \item \textbf{Redis:} Backend key-value storage engine for persistent data storage
    \item \textbf{Threading:} Python threading for background processes (gossip, anti-entropy, failure detection)
    \item \textbf{MD5 Hashing:} Consistent hashing implementation
\end{itemize}

\textbf{Key Modules:}
\begin{itemize}
    \item \texttt{worker.py}: Core node implementation with request handling and coordination logic
    \item \texttt{client.py}: Client interface for PUT/GET operations
    \item \texttt{HashRing.py}: Consistent hashing ring implementation
    \item \texttt{hinted\_handoff.py}: Temporary failure handling mechanism
    \item \texttt{merkle\_tree.py}: Anti-entropy synchronization
    \item \texttt{config.py}: System configuration (N, R, W parameters)
\end{itemize}

% Table I
\begin{table}[htbp]
\caption{A SUMMARY OF THE LIST OF TECHNIQUES OUR DYNAMOMINI USES AND THEIR RESPECTIVE ADVANTAGES.}
\begin{center}
\begin{tabular}{|p{0.25\linewidth}|p{0.3\linewidth}|p{0.3\linewidth}|}
\hline
\textbf{Problem} & \textbf{Technique} & \textbf{Advantage} \\
\hline
Partitioning & Consistent Hashing & Incremental Scalability \\
\hline
High Availability for writes & Vector clocks with reconciliation during reads & Version size is decoupled from update rates. \\
\hline
Handling temporary failures & Quorum-based algorithm \& using replicas & Provides high availability and durability guarantee when some of the replicas are not available. \\
\hline
Recovering from permanent failures & Using a Down Routing Table to maintain the set of down nodes & Synchronizes divergent replicas in the background. \\
\hline
Membership and failure detection & Gossip-based membership protocol and failure detection & Preserves symmetry and avoids having a centralized registry for storing membership and node liveness information. \\
\hline
\end{tabular}
\label{tab:summary}
\end{center}
\end{table} % [cite: 118, 119, 120]

\section{Implementation Challenges and Solutions}

Building a production-grade distributed storage system presents numerous challenges spanning consistency, availability, fault tolerance, and operational complexity. This section details the key challenges encountered during DynamoMINI development and the engineering solutions implemented.

\subsection{Challenge 1: Maintaining Routing Table Consistency}

\textbf{Problem:} In a decentralized system without a central authority, ensuring all nodes maintain a consistent view of cluster membership and topology is non-trivial. Inconsistent routing tables lead to requests routed to incorrect nodes, failed operations when nodes are unaware of new members, and split-brain scenarios during network partitions.

\textbf{Solution:} Implemented a gossip-based membership protocol where every node gossips with random peers periodically. Routing tables are tagged with vector clocks to detect and merge concurrent updates. Information propagates to all nodes in $O(\log N)$ rounds, achieving eventual consistency with typical convergence time under 30 seconds for clusters of 10+ nodes.

\subsection{Challenge 2: Handling Temporary vs. Permanent Failures}

\textbf{Problem:} Distinguishing between temporary failures (network hiccup, process restart) and permanent failures (hardware failure, node decommission) is critical. Premature removal causes unnecessary data movement; delayed removal impacts availability.

\textbf{Solution:} Implemented a two-tier failure detection system with a suspicion list for down nodes and active probing by background threads. Hinted handoff stores writes for temporarily unavailable nodes, automatically replaying upon recovery. Permanent removal requires explicit administrative action, preventing data loss from transient failures.

\subsection{Challenge 3: Achieving Eventual Consistency}

\textbf{Problem:} With asynchronous replication and network partitions, replicas inevitably diverge. Efficiently detecting and reconciling inconsistencies without full data scans is challenging.

\textbf{Solution:} Multi-layered consistency approach combining read repair (every GET triggers repair for stale replicas), Merkle tree anti-entropy (periodic background synchronization), vector clock reconciliation (automatic resolution for causally ordered updates), and timestamp-based tie-breaking for conflicts. Typical convergence time of 1-2 minutes for divergent replicas.

\subsection{Challenge 4: Dynamic Node Membership}

\textbf{Problem:} Adding or removing nodes requires careful coordination to avoid data loss and maintain availability. New nodes must acquire their key ranges; departing nodes must transfer data; preference lists must update atomically.

\textbf{Solution:} Implemented graceful scaling protocol where new nodes join via gossip, acquire data from successors after routing table convergence, and only begin accepting requests after full synchronization. Node removal involves pushing data to successors before departure notification. Achieved zero-downtime scaling operations with temporary 10-20\% latency increase during membership changes.

\subsection{Challenge 5: Quorum Coordination During Failures}

\textbf{Problem:} When multiple replicas in a preference list are simultaneously unavailable, achieving quorum becomes impossible, causing operation failures and violating availability guarantees.

\textbf{Solution:} Implemented sloppy quorum with hinted handoff. If W nodes in preference list are unavailable, writes go to first W healthy nodes in the ring. These temporary replicas store hints indicating true ownership, replaying upon original nodes' recovery. Maintained write availability even when up to N-1 nodes in a preference list are down.

\subsection{Challenge 6: Vector Clock Size Growth}

\textbf{Problem:} Vector clocks grow with the number of nodes that have updated a key. For highly contentious keys, vector clocks become unwieldy, consuming storage and bandwidth.

\textbf{Solution:} Implemented vector clock pruning with configurable maximum size (default: 10 entries). Prune least-recently-updated node entries when limit exceeded while retaining entries for nodes in current preference list. Bounded overhead with minimal impact on conflict detection accuracy.

\subsection{Challenge 7: Testing Distributed Behavior}

\textbf{Problem:} Testing distributed systems is difficult due to non-deterministic network behavior, race conditions, timing-dependent bugs, and difficulty reproducing specific failure scenarios.

\textbf{Solution:} Developed comprehensive testing infrastructure including network partition simulator, chaos testing (randomly killing nodes), consistency verification scripts, performance benchmarks under various failure scenarios, and visualization tools for hash ring and operation timelines. Identified and fixed numerous edge cases including race conditions in gossip protocol and routing table inconsistencies.

\section{Performance Analysis and Testing}

\subsection{Testing Methodology}
DynamoMINI was subjected to comprehensive testing to validate functionality, performance, and fault tolerance characteristics:

\textbf{Functional Testing:}
\begin{itemize}
    \item Basic PUT/GET operations across single and multiple nodes
    \item Replication verification—ensuring all N replicas receive writes
    \item Vector clock conflict detection and resolution
    \item Routing table consistency after gossip convergence
    \item Hinted handoff replay after node recovery
\end{itemize}

\textbf{Fault Injection Testing:}
\begin{itemize}
    \item Network partition simulation using \texttt{network\_partition.py}
    \item Random node failures during operations (chaos testing)
    \item Simultaneous multiple node failures
    \item Partition healing and consistency verification
\end{itemize}

\textbf{Performance Benchmarking:}
\begin{itemize}
    \item Latency measurements for PUT/GET under varying loads
    \item Throughput testing with concurrent clients
    \item Scalability testing—measuring performance as nodes increase
    \item Failure scenario impact on latency and availability
\end{itemize}

\subsection{Key Findings}
\textbf{Latency Characteristics:}
\begin{itemize}
    \item Median GET latency: 15-25ms (N=3, R=2)
    \item Median PUT latency: 20-30ms (N=3, W=2)
    \item P99 latency increases 3-5x during gossip convergence
\end{itemize}

\textbf{Availability:}
\begin{itemize}
    \item 99.9\%+ availability for reads/writes with up to 1 node failure
    \item Graceful degradation with multiple simultaneous failures
    \item Zero data loss observed during temporary failures with hinted handoff
\end{itemize}

\textbf{Consistency Convergence:}
\begin{itemize}
    \item Routing table convergence: $O(\log N)$ gossip rounds, typically 20-30 seconds
    \item Data consistency via read repair: 95\% of inconsistencies resolved within 1 minute
    \item Merkle tree synchronization: Scales efficiently to 10,000+ keys per node
\end{itemize}

\section{Use Cases and Applications}

\subsection{Amazon Shopping Cart}
The original motivation for Dynamo—Amazon's shopping cart service requires:
\begin{itemize}
    \item \textbf{High Write Availability:} Adding items must never fail, even during failures
    \item \textbf{Session Persistence:} Cart state must survive node failures
    \item \textbf{Eventual Consistency:} Acceptable to merge concurrent cart updates
\end{itemize}

DynamoMINI's design directly addresses these requirements through sloppy quorum (writes succeed even with node failures), replication (cart data survives failures), and vector clock reconciliation (merge concurrent cart modifications).

\subsection{Session Management}
Web applications requiring distributed session storage benefit from:
\begin{itemize}
    \item Low-latency session reads/writes (critical for response time)
    \item High availability (session loss degrades user experience)
    \item Partition tolerance (data centers may be isolated)
\end{itemize}

DynamoMINI provides sub-50ms operations, continues operating during partitions, and eventually reconciles session state across replicas.

\subsection{Distributed Configuration Store}
Applications needing replicated configuration data benefit from:
\begin{itemize}
    \item Eventual consistency (configuration updates can propagate gradually)
    \item High read availability (applications must always access configuration)
    \item Decentralized operation (no single point of failure)
\end{itemize}

\subsection{Social Media Interactions}
Counters for likes, views, follower counts where:
\begin{itemize}
    \item Exact counts less critical than availability
    \item High write throughput required
    \item Global distribution necessitates local writes
\end{itemize}

DynamoMINI's eventual consistency model and high write availability make it suitable for these use cases, accepting that counter values may temporarily differ across replicas.

\subsection{General Applicability}
DynamoMINI is well-suited for applications requiring:
\begin{itemize}
    \item \textbf{High Availability:} Service must remain operational during failures
    \item \textbf{Partition Tolerance:} System must handle network splits
    \item \textbf{Low Latency:} Sub-100ms response times desired
    \item \textbf{Scalability:} Linear scaling with nodes
    \item \textbf{Simple Data Model:} Key-value operations sufficient
\end{itemize}

Not recommended for applications requiring:
\begin{itemize}
    \item Strong consistency or ACID transactions
    \item Complex queries or joins
    \item Real-time consistency guarantees
\end{itemize}

\section{Related Work and Comparisons}

\subsection{Comparison with Other Distributed Systems}

\textbf{Cassandra:} Apache Cassandra is a production distributed database heavily inspired by Dynamo, adding features like CQL (query language), tunable consistency per-operation, and improved failure detection. DynamoMINI implements the core Dynamo principles demonstrated in Cassandra.

\textbf{Riak:} Another Dynamo-inspired system focusing on high availability and operational simplicity. Riak adds multi-data-center replication and improved conflict resolution (sibling management). DynamoMINI's vector clock implementation mirrors Riak's approach.

\textbf{DynamoDB (AWS):} Amazon's managed service evolved from the original Dynamo paper but with significant differences: strongly consistent reads option, automated partitioning management, and integration with AWS ecosystem. DynamoMINI stays closer to the original Dynamo design.

\textbf{Chord/Consistent Hashing:} DynamoMINI's ring-based partitioning is inspired by Chord DHT but simplified—Chord includes finger tables for $O(\log N)$ lookup, while DynamoMINI maintains full routing tables for $O(1)$ lookup at the cost of $O(N)$ memory.

\subsection{Trade-offs and Design Decisions}

\textbf{Eventual vs. Strong Consistency:} DynamoMINI sacrifices strong consistency for availability. Systems like Spanner or CockroachDB provide strong consistency using techniques like two-phase commit and synchronized clocks, but with higher latency and reduced availability during partitions.

\textbf{Full Routing Tables vs. Finger Tables:} DynamoMINI nodes maintain complete routing tables ($O(N)$ space) rather than Chord's finger tables ($O(\log N)$ space). This trades memory for lookup speed—beneficial for clusters of hundreds of nodes, but not for massive DHTs with millions of nodes.

\textbf{Gossip Convergence Time:} Gossip protocols provide eventual consistency with $O(\log N)$ convergence but no hard guarantees. Systems requiring immediate consistency must use consensus protocols (Paxos, Raft) at the cost of availability during partitions.

\section{Lessons Learned and Insights}

\subsection{Distributed Systems Complexity}
Building DynamoMINI revealed the inherent complexity of distributed systems:
\begin{itemize}
    \item \textbf{Testing Difficulty:} Race conditions and timing bugs are nearly impossible to reproduce deterministically
    \item \textbf{Failure Modes:} Partial failures (some nodes reachable, others not) are far more common than complete failures
    \item \textbf{Operational Overhead:} Even ``self-healing'' systems require careful monitoring and occasional manual intervention
\end{itemize}

\subsection{CAP Theorem in Practice}
The project validated CAP theorem predictions:
\begin{itemize}
    \item During network partitions, the system chose availability (accepting writes) over consistency
    \item Reconciliation complexity was significant—conflict resolution pushed to reads or application logic
    \item Tunable consistency (N, R, W parameters) allowed trading latency for consistency on per-application basis
\end{itemize}

\subsection{Design Principles for Distributed Storage}
\begin{itemize}
    \item \textbf{Decentralization:} Symmetric peer-to-peer design eliminates single points of failure and simplifies operation
    \item \textbf{Incremental Scalability:} Consistent hashing enables adding/removing nodes with minimal data movement
    \item \textbf{Heterogeneity:} Virtual nodes allow balancing load across nodes with different capacities
    \item \textbf{Optimistic Replication:} Accepting temporary inconsistency dramatically improves write availability
\end{itemize}

\section{Future Enhancements}

While DynamoMINI implements core Dynamo functionality, several enhancements would improve production readiness:

\subsection{Performance Optimizations}
\begin{itemize}
    \item \textbf{Request Pipelining:} Batch multiple operations to same node to reduce RPC overhead
    \item \textbf{Parallel Reads:} Issue read requests to all N replicas simultaneously, using first R responses
    \item \textbf{Connection Pooling:} Reuse persistent connections rather than establishing per-request
    \item \textbf{Bloom Filters:} Reduce storage lookups for non-existent keys
\end{itemize}

\subsection{Operational Enhancements}
\begin{itemize}
    \item \textbf{Metrics and Monitoring:} Expose Prometheus-compatible metrics for latency, throughput, error rates
    \item \textbf{Administrative Interface:} Web UI for viewing ring state, node health, replication status
    \item \textbf{Automated Capacity Planning:} Detect hotspots and recommend virtual node adjustments
    \item \textbf{Backup and Restore:} Snapshot mechanisms for disaster recovery
\end{itemize}

\subsection{Advanced Features}
\begin{itemize}
    \item \textbf{Range Queries:} Support for ordered key traversal (requires ordered partitioning scheme)
    \item \textbf{Multi-Data-Center Replication:} Topology-aware placement for geographic distribution
    \item \textbf{Compression:} Reduce storage and network overhead for large values
    \item \textbf{Client Libraries:} Production-ready SDKs in multiple languages with retry logic and connection pooling
\end{itemize}

\subsection{Consistency Improvements}
\begin{itemize}
    \item \textbf{Last-Write-Wins Optimization:} For applications not needing conflict detection, simple timestamp comparison suffices
    \item \textbf{CRDT Support:} Conflict-free replicated data types for automatic reconciliation of specific data structures
    \item \textbf{Causal Consistency:} Maintain causality guarantees for related operations
\end{itemize}

\section{Conclusion}
DynamoMINI successfully demonstrates the principles and challenges of building a highly available distributed storage system inspired by Amazon's Dynamo. Through implementation of consistent hashing, vector clocks, quorum-based replication, gossip protocols, hinted handoff, and Merkle tree anti-entropy, the system achieves the core goals of high availability, partition tolerance, and eventual consistency.

The project validates the feasibility of Dynamo's design choices while revealing the inherent complexity of distributed systems. Key achievements include:

\begin{itemize}
    \item \textbf{Decentralized Architecture:} Peer-to-peer design with no single point of failure
    \item \textbf{High Availability:} Operations succeed even during node failures and network partitions
    \item \textbf{Scalability:} Linear scalability through consistent hashing and symmetric design
    \item \textbf{Fault Tolerance:} Automatic failure detection, hinted handoff, and background repair
    \item \textbf{Tunable Consistency:} Configurable N, R, W parameters for application-specific trade-offs
\end{itemize}

The experience of building DynamoMINI underscores the sophisticated engineering required for production distributed systems. While the core algorithms are well-understood, their implementation demands careful attention to edge cases, failure modes, and performance characteristics. The system serves both as an educational tool for understanding distributed storage principles and as a foundation for more sophisticated implementations.

DynamoDB's influence on modern distributed databases is evident in systems like Cassandra, Riak, and Voldemort. DynamoMINI's implementation confirms that these principles—while complex—are achievable and provide a solid foundation for building scalable, available storage systems. The trade-offs made (eventual consistency, conflict resolution complexity) are well-suited for many real-world applications where availability trumps immediate consistency.

Future work could extend DynamoMINI toward production readiness through performance optimizations, operational tooling, and advanced features. However, the current implementation successfully demonstrates that the core Dynamo architecture is both practical and effective for building distributed storage systems that prioritize availability and partition tolerance.

\section{Acknowledgement}
We are grateful for the opportunity to this project as part of a Distributed Systems course taken by Dr. Kishore Kothapalli. % [cite: 157, 158]

\begin{thebibliography}{00}
\bibitem{b1} Giuseppe DeCandia et al., ``Dynamo: Amazon's highly available key-value store.''
\bibitem{b2} Amazon DynamoDB: A Scalable, Predictably Performant, and Fully Managed NoSQL Database Service.
\end{thebibliography}

\end{document}